\label{sec:metodologia}

\section{Fluxo da informação}\label{sec:fluxo}

A Figura~\ref{fig:fluxo} descreve o fluxo proposto e define os dois pontos de
intervenção. As instruções do desenvolvedor e a requisição do usuário chegam ao filtro por
\textbf{canais declaradamente distintos}: o canal de instrução carrega a
política de classificação, e o canal de dado carrega o conteúdo não confiável,
delimitado por uma tag explícita. O filtro emite uma decisão e, quando libera a
requisição, também a lista dos trechos que julgou perigosos. Esses trechos são
marcados tipograficamente no texto original, que segue íntegro, sem qualquer
supressão, e repassados ao componente principal, o qual recebe, junto ao
texto anotado, a instrução de tratar trechos marcados como dado de risco. O
desfecho é medido na resposta do componente principal.

\begin{figure}[htb]
\centering
\resizebox{\textwidth}{!}{% Fluxo da informação: filtro com canais separados -> marcação seletiva ->
% componente principal -> ponto de medição. Incluído por TG3.tex via \input.
\begin{tikzpicture}[
  font=\footnotesize,
  node distance=6mm,
  caixa/.style={draw, rounded corners=2pt, align=center, inner sep=3pt,
                minimum height=7mm, text width=21mm},
  canal/.style={draw, dashed, align=center, inner sep=2.5pt,
                minimum height=6mm, text width=25mm},
  grupo/.style={draw, thick, rounded corners=3pt, inner sep=4pt},
  medida/.style={draw, thick, rounded corners=2pt, align=center, inner sep=3pt,
                 minimum height=7mm, text width=25mm},
  seta/.style={-latex, thick},
]

% ---- Entradas -------------------------------------------------------------
\node[caixa] (conf)   {Instruções do desenvolvedor\\(confiável)};
\node[caixa, below=4mm of conf] (naoconf) {Requisição do usuário\\(não confiável)};

% ---- Filtro: dois canais separados ---------------------------------------
\node[canal, right=12mm of conf, yshift=0mm] (cinstr)
      {\textbf{canal de instrução}\\política de classificação};
\node[canal, below=3mm of cinstr] (cdado)
      {\textbf{canal de dado}\\delimitado em \texttt{<data>}};
\node[grupo, fit=(cinstr)(cdado), label={[font=\scriptsize\bfseries]above:Filtro de entrada (modelo aberto)}] (filtro) {};

% ---- Saída do filtro ------------------------------------------------------
\node[caixa, right=11mm of filtro, text width=19mm] (decisao)
      {Decisão\\\emph{bloquear} / \emph{passar}};
\node[caixa, right=7mm of decisao, text width=23mm] (marca)
      {Marcação seletiva\\dos trechos suspeitos (\texttt{\^{}})};

% ---- Componente principal e medição --------------------------------------
\node[caixa, right=7mm of marca, text width=21mm] (principal)
      {Componente principal\\(\emph{backend} LLM)};
\node[medida, right=7mm of principal] (medida)
      {Medição do desfecho\\ASR, utilidade,\\FN, FP};

% ---- Arestas --------------------------------------------------------------
\draw[seta] (conf.east)    -- ++(4mm,0) |- (cinstr.west);
\draw[seta] (naoconf.east) -- ++(4mm,0) |- (cdado.west);
\draw[seta] (filtro.east)  -- (decisao.west);
\draw[seta] (decisao.east) -- node[above, font=\scriptsize] {passa} (marca.west);
\draw[seta] (marca.east)   -- (principal.west);
\draw[seta] (principal.east) -- node[above, font=\scriptsize] {resposta} (medida.west);

% bloqueio: encerra o fluxo
\node[below=7mm of decisao, font=\scriptsize] (fim) {fim do fluxo};
\draw[seta] (decisao.south) -- node[right, font=\scriptsize] {bloqueia} (fim.north);
\draw[seta] (fim.east) -| ([xshift=-3mm]medida.south);

% instrução de leitura da marcação
\draw[seta, dotted] (marca.north) -- ++(0,5mm)
      -| node[pos=0.25, above, font=\scriptsize] {instrução de cautela} (principal.north);

\end{tikzpicture}
}
\caption{Fluxo da informação. O filtro recebe instrução e dado por canais
separados; ao liberar a requisição, marca os trechos suspeitos e os repassa
anotados ao componente principal. A medição ocorre na saída do sistema, não na
classificação do filtro.}
\label{fig:fluxo}
\end{figure}

O Apêndice~\ref{ape:exemplos} percorre esse fluxo sobre dois itens reais, um
atacado e um benigno, com o texto literal do que cada estágio recebeu e
respondeu sob cada configuração.

\section{Primeiro fator: delimitação na entrada do filtro}\label{sec:fator1}

O primeiro fator controla como o conteúdo não confiável é apresentado ao
filtro, mantendo constante a política de classificação. Sem delimitação, o
conteúdo é concatenado à requisição, como faz qualquer \textit{guardrail}
convencional. Com delimitação, ele é envolvido em uma marcação estrutural
explícita, isto é, a tag XML \texttt{<data>}, enquanto a política de classificação
permanece fora de qualquer tag, e o \textit{prompt} do filtro declara que o
conteúdo interno à tag é dado a ser inspecionado, nunca instrução a ser seguida.

Um requisito de implementação decide se este fator mede o que pretende: o
conteúdo não confiável precisa ter seus caracteres \texttt{<} e \texttt{>}
escapados antes de entrar na tag. Sem isso, bastaria ao atacante incluir um
\texttt{</data>} literal para encerrar a delimitação e voltar a ser lido como
instrução. Nesta situação, o experimento estaria medindo uma delimitação que não delimita,
situação pior do que a ausência de delimitação, por conferir falsa confiança ao
resultado.

\section{Segundo fator: marcação seletiva na saída do filtro}\label{sec:fator2}

O segundo fator controla o que o filtro repassa ao componente principal quando
decide \textbf{liberar} a requisição. Sem marcação, ele se comporta como o
\textit{guardrail} clássico: decide passar ou bloquear e não altera o texto. Com
marcação, o filtro marca tipograficamente, com um caractere intercalado, apenas
os trechos que ele próprio apontou como perigosos, e o \textit{prompt} do
componente principal declara explicitamente que trechos marcados representam
risco e devem ser tratados como dado, jamais como instrução.

É neste fator que reside a proposta inovadora do trabalho, e a diferença em relação às
defesas por sanitização está na decisão sobre o trecho suspeito: elas o
\textbf{removem} \cite{shi2025promptarmor, wang2025datafilter, geng2025pisanitizer};
aqui ele é \textbf{preservado e sinalizado}. Operacionalmente, o contrato de
saída do filtro inclui a lista literal dos trechos suspeitos, e a marcação é
aplicada ao texto original por casamento aproximado, a mesma técnica de
localização que o PromptArmor emprega para remover \cite{shi2025promptarmor},
aqui empregada para preservar.

A marcação é \textbf{seletiva}, e não integral, por duas razões. Marcar o bloco
inteiro não discrimina: se tudo está marcado, nada se destaca, e a marca degrada
a leitura do conteúdo legítimo. Em contrapartida, a marcação seletiva só protege
aquilo que o filtro efetivamente identificou. Quando ele não aponta trecho
algum, a condição degrada-se graciosamente para o \textit{guardrail} clássico.
Essa dependência da qualidade da identificação é uma propriedade do método, e
não um defeito da implementação; medi-la é parte do que o experimento faz.

Com os dois fatores definidos, o cruzamento entre eles produz um desenho
matricial $2 \times 2$: quatro condições, uma para cada combinação de presença e
ausência da delimitação na entrada e da marcação seletiva na saída
(Tabela~\ref{tab:condicoes}). Os fatores são binários e independentes, e as
quatro condições são executadas sobre os mesmos itens, o que permite isolar o
efeito de cada intervenção e o da combinação das duas. A célula superior
esquerda, o \textit{prompt} simples, é o \textit{guardrail} clássico: o mesmo
filtro, com a mesma política, sem separação de canal em nenhum dos dois pontos.
É contra ela que se mede o efeito da separação. A referência sem filtro, que
não tem \textit{guardrail} algum, não pertence à tabela e é descrita na
Seção~\ref{sec:linhas}.

\begin{table}[H]
\centering
\caption{As quatro condições do desenho fatorial.}
\label{tab:condicoes}
\begin{tabular}{lll}
\toprule
 & Sem marcação & Com marcação seletiva \\
\midrule
Sem delimitação & \textit{prompt} simples & marcação \\
Com delimitação & delimitação & delimitação + marcação \\
\bottomrule
\end{tabular}
\vspace{2pt}
\begin{minipage}{\textwidth}
\footnotesize\raggedright
\textbf{Legenda.} Linhas: primeiro fator, a delimitação na entrada do filtro
(Seção~\ref{sec:fator1}). Colunas: segundo fator, a marcação seletiva na saída
do filtro (Seção~\ref{sec:fator2}). Cada célula nomeia a condição como ela
aparece nas figuras e tabelas de resultados. A composição exata das mensagens
de cada condição está na Tabela~\ref{tab:ape_montagem}, no
Apêndice~\ref{ape:exemplos}.
\end{minipage}
\end{table}

\section{Linhas: modelos e referências}\label{sec:linhas}

Cada bateria é executada sobre dois modelos abertos, e o modelo instancia o
\textbf{sistema inteiro}, tanto o filtro quanto o componente principal, de
modo que cada bateria descreva um sistema homogêneo. Os modelos são o
Llama-3.1-8B-Instruct e o gpt-oss-20b, este último uma arquitetura de mistura de especialistas com 21
bilhões de parâmetros totais e 3,6 bilhões ativos por \textit{token}.

Somam-se as referências. A primeira é a condição \textbf{sem filtro}, que
estabelece a taxa de sucesso de ataque sem defesa alguma; sem ela não haveria
denominador contra o qual expressar qualquer ganho. A segunda é o
\textbf{filtro fine-tunado} da família gpt-oss: o gpt-oss-safeguard, treinado
para classificação de segurança sobre exatamente o mesmo \textit{backbone} do
gpt-oss-20b, entra no papel de filtro nas mesmas quatro condições da
Tabela~\ref{tab:condicoes}, com o gpt-oss-20b como componente principal. Assim
a comparação isola, célula a célula, o que o treinamento disponível na família
acrescenta e o que a separação de canal ainda acrescenta sobre um filtro
treinado, e o faz por medição direta, sem depender de números publicados por
terceiros. O gpt-oss-safeguard é um classificador de segurança de conteúdo
guiado pela política do \textit{prompt} de sistema, mas não é um detector treinado
especificamente para \textit{prompt injection}, portanto, ele recebe a mesma política que
os modelos abertos.

A família Llama não recebe referência equivalente, e a razão é de
disponibilidade. Os detectores treinados especificamente para \textit{prompt
injection} sobre esse \textit{backbone} existem como pesos publicados
\cite{wang2025datafilter, geng2025pisanitizer, jacob2025promptshield}, mas
nenhum é servido por API: usá-los exigiria hospedar um modelo de oito bilhões
de parâmetros em infraestrutura própria, o que sairia do custo e da
reprodutibilidade que o trabalho adota. Os detectores pequenos o bastante para
rodar sem servidor dedicado são codificadores treinados à parte, e não modelos
da família avaliada. Essas defesas permanecem, portanto, como referência de
literatura, discutidas no Capítulo~\ref{cap:fundamentacao}, e a comparação
medida é a que a família gpt-oss permite. O \textit{guard} de conteúdo da
família Llama, o Llama Guard~4, não entra no desenho por ser derivado do
Llama~4 Scout, e não do Llama~3.1.

O desenho completo compreende, portanto, \textbf{quatorze execuções} por
benchmark: as quatro condições sobre os dois modelos abertos, duas referências
sem filtro e as quatro condições com o filtro fine-tunado.

\section{Benchmarks}\label{sec:benchmarks}

A avaliação usa três benchmarks públicos e citáveis, cada um responsável por um
eixo distinto do desfecho:

\begin{itemize}
  \item \textbf{BIPIA} \cite{yi2025bipia}: injeção indireta em cinco tarefas
  de texto (correio eletrônico, perguntas sobre página web, tabelas, sumarização
  e código). Fornece simultaneamente a taxa de sucesso do ataque e a utilidade na
  tarefa legítima, e é o único dos três que exercita o fluxo completo até a
  resposta do componente principal. É o benchmark primário.
  \item \textbf{SEP} \cite{zverev2025separation}: mede diretamente o grau de
  separação entre instrução e dado, comparando o comportamento do modelo quando
  a mesma sentença ocupa a posição de instrução e a de dado. É o instrumento que
  liga o resultado empírico à propriedade formal que a intervenção pretende
  melhorar.
  \item \textbf{NotInject} \cite{li2024injecguard}: conjunto exclusivamente
  benigno, construído com palavras-gatilho que induzem detectores ao erro. Mede
  o sobre-bloqueio, isto é, o preço pago em falsos positivos por qualquer ganho
  de detecção.
\end{itemize}

A amostragem é estratificada por rótulo com semente fixa, e as suítes são
congeladas antes da execução.

\section{Métricas e protocolo}\label{sec:protocolo}

Registram-se métricas em dois pontos. \textbf{No filtro}: taxa de detecção a
1\% de falsos positivos, sobre-bloqueio sobre o NotInject, custo em
\textit{tokens} e tempo de geração p50/p95, informado pelo provedor de cada
chamada, além da \textbf{taxa de saída inválida}: modelos
pequenos frequentemente violam o formato de saída pedido, e itens
inválidos são contabilizados como não bloqueio.
\textbf{Na saída do componente principal}: taxa de sucesso do ataque, utilidade
preservada na tarefa benigna, e a decomposição dos erros em falso-negativo
(ataque liberado e executado) e falso-positivo (requisição benigna bloqueada ou
degradada pela marcação). A resposta vazia do componente principal, que ocorre
quando um modelo de raciocínio esgota o teto de \textit{tokens} antes de
responder, é tratada como saída inválida, e não como bloqueio: não executa o
ataque, não cumpre a tarefa e não entra no sobre-bloqueio, que mede apenas a
decisão do filtro. Sua frequência é reportada à parte, como a da saída
inválida do filtro. A tarefa conta como cumprida quando ao menos metade
das palavras de conteúdo da resposta de referência aparece, como palavra
inteira, na resposta produzida; referências curtas, como um número, precisam
aparecer por inteiro. Palavra, aqui, é uma sequência de letras e dígitos, de
modo que pontuação e marcação de formatação coladas a ela não impedem o
casamento. A regra evita dois erros de um critério mais simples: separar as
palavras só por espaço reprovaria respostas corretas escritas em negrito ou
seguidas de ponto, e procurar a referência como trecho do texto aprovaria
``4'' dentro de qualquer número que contivesse esse dígito.

Explicando a taxa de detecção a 1\% de falsos positivos: a forma de medição
vem de \citeonline{jacob2025promptshield}, que argumentam que, no tráfego
real, as entradas benignas superam de longe os ataques, de modo que um
detector só é implantável se quase nunca bloquear conteúdo legítimo. Em vez
de julgar o detector pelo limiar padrão da classificação, eles calibram o
limiar para uma taxa alvo de falsos positivos (1\%, 0{,}5\%, 0{,}1\% e
0{,}05\%) e reportam a detecção obtida em cada uma. Aqui adota-se o alvo de
1\%, o menos exigente da escala, porque as amostras têm 150 itens benignos e
alvos menores não teriam resolução. Além da decisão, o filtro atribui a cada
item um escore de risco entre 0 e 1. Ordenados os itens por escore,
pergunta-se a partir de que valor se deveria bloquear para que, no máximo,
1\% dos itens benignos fossem bloqueados, e reporta-se a fração dos atacados
que esse limiar captura. Um exemplo em miniatura: dez itens, cinco
benignos com escores 0{,}0, 0{,}1, 0{,}2, 0{,}3 e 0{,}5, e cinco atacados com
escores 0{,}4, 0{,}6, 0{,}7, 0{,}8 e 0{,}9. Bloqueando a partir de 0{,}6,
caem quatro atacados e nenhum benigno: detecção de 0{,}80 com 0\% de falsos
positivos. Baixando o limiar para 0{,}5, o benigno de escore 0{,}5 passa a
ser bloqueado, e a taxa de falsos positivos salta para 20\%, fora do
orçamento. O maior limiar que respeita a restrição é, portanto, 0{,}6, e a
métrica vale 0{,}80; o atacado de escore 0{,}4 fica sem detecção, e esse é o
preço da restrição. Com os 150 benignos do BIPIA, 1\% equivale a no máximo
um benigno bloqueado. A métrica é, assim, uma reanálise dos escores após a
execução: o limiar que o filtro de fato aplicou na bateria é o que alimenta
as demais métricas, enquanto esta mede quanto o classificador detectaria se o
limiar fosse ajustado sob a restrição. Ela depende da resolução do escore: um
filtro que só emite três valores distintos raramente encontra um limiar
intermediário, e o número reportado cai a zero, convenção que também segue
a referência: quando nenhum limiar atinge a taxa alvo, a detecção é dada
como nula.

O efeito é exibido através da \textbf{diferença} contra duas referências distintas: 
sem filtro e \textit{prompt} simples, porque elas respondem a perguntas diferentes. 
Contra a condição \textbf{sem filtro}, obtém-se o quanto o \textit{guardrail} protege no total; contra o
\textit{prompt} \textbf{simples} (que já é um \textit{guardrail} completo)
obtém-se o que é atribuível especificamente à separação de canal, que é a
pergunta de pesquisa.

A bateria inteira é executada \textbf{dez vezes}, com a mesma configuração e
sobre os mesmos itens. A repetição existe porque o modelo não é
determinístico, nem mesmo com temperatura zero: o \textit{prompt} simples e a
marcação enviam ao filtro exatamente o mesmo texto, já que a marcação só age
depois da decisão, e ainda assim a decisão sobre um mesmo item muda de uma
execução para outra. O desfecho de cada item passa a ser a média das suas dez
repetições, isto é, a fração delas em que o ataque foi executado ou a tarefa
foi cumprida, e é sobre essas médias que todas as grandezas são calculadas.

Os intervalos de confiança são de 95\% e vêm do \textit{bootstrap} BCa
\cite{efron1987better}, com 10.000 reamostragens dos itens, tanto para as
taxas absolutas quanto para as diferenças entre condições e para os efeitos
do desenho fatorial. Como todas as condições são avaliadas sobre exatamente
os mesmos itens, a diferença entre duas condições é calculada item a item, e
cada reamostragem sorteia itens levando junto os seus valores nas condições
comparadas. O intervalo é, assim, pareado: a variação que vem do item, e não
da condição, sai da conta. A unidade sorteada é o item, e não a execução.
Cada item entra com a média das suas repetições, e a variação entre
repetições viaja com ele; contar cada execução como um item independente
trataria dez medições do mesmo item como dez itens e estreitaria o intervalo
sem fundamento. O BCa corrige o viés e a assimetria da distribuição
reamostrada. Quando menos de cinco itens se afastam do valor mais comum, como
numa taxa de sobre-bloqueio quase sempre nula, a reamostragem não tem o que
sortear e devolveria um intervalo de largura nula ou quase, que afirmaria uma
precisão que o dado não tem. Nesses casos não se reporta intervalo: a barra
aparece sem haste, e o valor se lê como a contagem que ele é, por exemplo um
bloqueio em 1.500 decisões. O valor central é sempre a média sobre os itens, que
coincide com a taxa agregada das dez repetições. Nas figuras, o intervalo
aparece como a haste vertical em torno de cada barra ou ponto, e nas
diferenças a leitura relevante é se a haste contém ou não o zero: quando não
contém, a diferença não se explica pelo acaso, seja o da amostragem dos itens,
seja o da variação entre execuções.

O desenho fatorial permite, além das comparações contra o \textit{prompt}
simples, estimar o efeito de cada fator usando as quatro condições de uma vez.
O \textbf{efeito principal} da delimitação é a média das duas condições com
delimitação menos a média das duas sem; o da marcação é definido da mesma
forma. A \textbf{interação} é a diferença entre o ganho da marcação com e sem
delimitação, e vale zero quando os dois fatores simplesmente se somam. Cada
efeito é calculado item a item e resumido pela média; como cada item entra
nas quatro condições, o intervalo fica mais estreito do que o de qualquer
comparação entre duas condições. O intervalo é o mesmo BCa sobre os itens. O
Apêndice~\ref{ape:incerteza} desenvolve as contas dos intervalos sobre
exemplos numéricos e mostra como as dez repetições entram nelas.

A semente da amostragem é fixa, de modo que as dez repetições usam os mesmos
itens, e todos os registros por item são persistidos, o que garante a
reprodutibilidade da análise. O registro de cada item inclui o conteúdo não
confiável como cada estágio o recebeu: delimitado, na entrada do
filtro; anotado, na entrada do componente principal, além dos trechos
efetivamente marcados, de modo que a transformação aplicada pelo
\textit{guardrail} possa ser auditada após a execução, e não apenas inferida da
condição.

\section{Situando o resultado frente ao \textit{fine-tuning}}\label{sec:ancora}

Saber o quanto a configuração recupera exige compará-la a um modelo treinado
para o papel de filtro, e a comparação é feita por medição direta, na mesma
bateria, sobre os mesmos itens e com as mesmas métricas. O gpt-oss-safeguard
compartilha o \textit{backbone} do gpt-oss-20b e roda as mesmas quatro
condições, o que permite ler a diferença em dois sentidos: entre as duas
linhas, célula a célula, está o que o treinamento de segurança acrescenta;
dentro da linha do safeguard está o que a separação de canal acrescenta a um
filtro já treinado. A segunda leitura é a que interessa à pergunta de
pesquisa, porque responde se a configuração ainda tem efeito quando o filtro
não é um modelo de uso geral.

Resultados publicados por terceiros sobre os mesmos benchmarks abertos
\cite{chen2025metasecalign, wang2025datafilter} entram apenas como contexto
secundário, e cada número externo é registrado com benchmark, métrica, valor e
fonte, de modo que sua procedência acompanhe o resultado.

\section{Execução e consolidação}\label{sec:analise}

A bateria foi executada integralmente sobre os três benchmarks, com quatorze
execuções por benchmark: as quatro condições da Tabela~\ref{tab:condicoes} nos
dois modelos abertos, as duas referências sem filtro e as quatro condições com
o filtro fine-tunado, ou seja, 42 execuções por repetição. Com as dez
repetições, a bateria soma \textbf{420 execuções e 131.460 registros}, um por
item em cada repetição. O BIPIA e o SEP foram avaliados sobre 300 itens
amostrados de forma estratificada; do NotInject, que é menor, foram usados os
339 itens completos.

Na consolidação dos registros são calculados todos os intervalos de confiança, cada condição
contra as duas referências descritas na Seção~\ref{sec:protocolo}, para taxa
de sucesso do ataque e para utilidade, e não apenas os de um recorte.

\subsection{Orçamento de geração}

Toda chamada de modelo recebe um teto de geração: 2.048 \textit{tokens} no
filtro e 8.192 no componente principal, iguais para todas as condições. A
finalidade do teto é limitar custo e latência por item e garantir que cada
chamada termine; o valor, porém, não é indiferente ao resultado. Um dos modelos
abertos adotados é modelo de raciocínio: ele consome parte do orçamento de
geração antes de emitir a resposta e, se o teto é alcançado antes do fim, a
interface devolve conteúdo vazio, acompanhado da indicação de truncamento. Uma
resposta vazia não contém a testemunha do ataque, e por isso conta como ataque
não executado; tampouco cumpre a tarefa. O efeito não é uniforme: o
esvaziamento cresce com o comprimento do \textit{prompt}, e as condições que
acrescentam texto, justamente as que implementam a separação de canal, seriam
as mais truncadas, produzindo um viés na direção da hipótese.

O teto de 8.192 \textit{tokens} no componente principal foi escolhido para
que a falta de orçamento não decida o resultado: a esse valor, o esvaziamento
fica restrito a uma pequena fração das respostas, em que o modelo esgota o teto
raciocinando. Essas respostas vazias contam como saída inválida do
componente, com frequência reportada à parte (Seção~\ref{sec:res_validade}).
No filtro, o teto de 2.048 \textit{tokens} basta para a decisão em JSON; as
chamadas que o atingem sem decisão legível contam como não bloqueio, a regra
conservadora adotada para toda saída inválida. O registro de cada item guarda
o motivo de término dos dois estágios.
